\scanpage{5}
of the Artin ring $B_n=H^{\unsure{o|0}}(X_n,\underline{O}_{X_n})$. Assuming $U$ hence $X$ hence $X_n$ to be connected, this Artin ring is local, and it's Brauer group coincides with the Brauer group of its residue field $k_n$. Of course, \deleted{thus} for an element of $P(k)=\lim P_n(k)$, these vari\deleted{\unsure{e}}ous obstructions match together and come from an element of the Brauer group of the residue fi\unsure{el|le}d $k'=\handwritten{\bigcap_n}k_n$ of the ring $H^o(X,\underline{O}_X)=\underset{\handwritten{n}}{\underset{\handwritten{\longleftarrow}}{\lim}}H^o(X_n,\underline{O}_{X_n})$, \deleted{\unsure{of}} which is just the normalisation of $A$. If for instance $A$ is \uline{normal}, \inserted{\handwritten{then}} $k'=k$ and the obstruction lies in $\operatorname{Br}(k)$.

Whatever this be, defining $\unsure{\mathbb{Q}|Q}$ as before, we still do get an injective homomorphism (6), but even if (5) is injective, we cannot be sure that (6) is bijective. We are however if $k$ \unsure{is|ss} separable closed. Therefore, we still \deleted{get} can more or less describe $\deleted{\unsure{P}}\handwritten{Q}$ (or at least $\deleted{\unsure{P}}\handwritten{Q}(k)$) in termes of local Picard groups, in terms of \inserted{Galois} descent from the completion $A'$ of the strict henselization of $A$: if $U'=S'-\handwritten{\{}s'\handwritten{\}}$, where $S'=\operatorname{Spec}(\unsure{A'|A(})$, and if
\[
\handwritten{\pi}=\operatorname{Gal}(\bar{k}/k)\ ,
\tag*{(9)}
\]
we will have an isomorphism
\[
\deleted{\operatorname{Pic}}\mathbb{Q}(k)\handwritten{\simeq}\underline{\operatorname{Pic}}(U')^{\handwritten{\pi}}\ ,
\tag*{(10)}
\]
\deleted{\unsure{This is perf} provided $D'$ is torsion free, for instance (5) injective,}

This is perfectly natural and in accordance with the familiar phenomena when dealing with \uline{global} Picard schemes. Therefore $Q$ looks like \unsure{the|ahe} natural object to be called a \deleted{Picard} local Picard scheme.

To make sure it makes a sense, one will have however to analyze somewhat the inverse system $(P_n)$, and the morphism (5). Moreover, to feel really secure, one should check that up to canonical isomorphism, the group scheme $Q$ together with (10) does not depend on the choice of
% archive sheet number: 4
